\documentclass[10pt,a4paper]{amsart}
\usepackage[margin=19mm]{geometry}
\usepackage{amsmath,amssymb,mathtools}
\usepackage[scr=rsfs,cal=euler]{mathalfa}
\usepackage{xfrac}
\usepackage[unicode,hidelinks,bookmarks=false]{hyperref}
\newcommand{\ie}{i.e.\ }
\newcommand{\cf}{cf.\ }
\newcommand{\resp}{respectively\ }
\newcommand{\Subsection}{\subsection}
\providecommand{\nobreakdash}{\nobreak\hbox{-}}
\setlength{\emergencystretch}{2em}
\multlinegap0pt
\usepackage{bm}
\usepackage{bbm}
\usepackage{dsfont}
\usepackage{xspace}
\usepackage{cancel}
\usepackage{mathtools}
\usepackage{amsthm}
\makeatletter
\@ifundefined{theorem}{\newtheorem{theorem}{Theorem}}{}
\@ifundefined{lemma}{\newtheorem{lemma}[theorem]{Lemma}}{}
\makeatother
\newcommand{\Ent}{\textnormal{Ent}}
\newcommand{\E}{\mathbb{E}}
\newcommand{\bZ}{\boldsymbol{Z}}
\newcommand{\bfA}{\mathbf{A}}
\newcommand{\bz}{\boldsymbol{z}}
\newcommand{\diff}{\textnormal{d}}
\title[Dimension-free Bounds for Sum of Dependent Matrices and Oper]{Dimension-free Bounds for Sum of Dependent Matrices and Operators with Heavy-Tailed Distribution}
\author{Shogo Nakakita, Pierre Alquier, Masaaki Imaizumi}
\date{Source: arXiv:2210.09756v3 (CC BY 4.0)}
\begin{document}
\maketitle

\noindent\emph{Source and licence.} Dimension-free Bounds for Sum of Dependent Matrices and Operators with Heavy-Tailed Distribution. Version arXiv:2210.09756v3 (CC BY 4.0), distributed under the Creative Commons Attribution 4.0 International licence, as recorded in the arXiv licence metadata for this exact version and shown on the abstract page https://arxiv.org/abs/2210.09756v3. Full rights notice: licenses/arxiv-2210.09756-CC-BY-4.0.txt.\par\smallskip

\noindent\textit{Editorial scope.} Counted unit: the Lemma whose source label is \texttt{lemma:expmoment} in the article (source0.tex lines 661--667, source label \texttt{lemma:expmoment}), delivered together with the article's own complete original proof of it (source0.tex lines 669--759, one \texttt{proof} environment of 91 lines). No mathematics is written, repaired or re-derived: statement and proof are the article's own text line for line. Required context retained uncounted: none. Every definition, display and label of those lines is retained unchanged. Omitted from this package and not redistributed: 1--660, 668--668, 760--2671. Nothing inside a retained excerpt is omitted or altered: every retained body is delivered exactly as the article's lines are written, so no omission receipt of any kind accompanies this package. Citations: the retained excerpts carry no citation command at all, so no bibliography entry of the article is transcribed and no reference is redistributed. Reference adaptation: none was needed and none was made; every reference inside the delivered unit resolves inside it, so no unresolved reference or citation is produced. Printed slips kept literal: no printed word, symbol, hypothesis or number of the counted statement or of its proof was altered. Layout and class: the article's own class is \texttt{article} with packages including geometry, titling, amsfonts, amsmath, amsthm, mathrsfs, stmaryrd, mathtools, natbib, url, tikz, color, newtxtext, newtxmath; the wrapper is the lane's pinned \texttt{amsart} editorial wrapper at 10pt on A4, which restates the theorem-like environments the retained text needs and adds the article's own macro definitions that the retained text needs (\texttt{\textbackslash E}, \texttt{\textbackslash Ent}, \texttt{\textbackslash bZ}, \texttt{\textbackslash bfA}, \texttt{\textbackslash bz}, \texttt{\textbackslash diff}), with extra wrapper packages (bm, bbm, dsfont, xspace, cancel, mathtools). The delivered numbering is the unit's own. Assets: no retained span contains a figure environment, a graphics-inclusion command, a PDF-page inclusion, a table or a TikZ picture; no image file is copied into this package, the package declares no asset dependency and no image file is redistributed with it. Licence: version arXiv:2210.09756v3 of this article is distributed under the Creative Commons Attribution 4.0 International licence, as recorded in the arXiv licence metadata for this exact version and shown on the abstract page https://arxiv.org/abs/2210.09756v3. A full rights notice is delivered at licenses/arxiv-2210.09756-CC-BY-4.0.txt. Characters: every retained excerpt is pure ASCII, so the delivered engine, XeLaTeX with its default Latin Modern text font, reports no missing character.
\begin{lemma}[exponential integrability of quadratic forms]\label{lemma:expmoment}
    Suppose that $\bfA$ is a $d\times d$ positive semi-definite matrix and $\bZ$ is a $d$-dimensional random vector whose distribution $\nu$ satisfies a log-Sobolev inequality with constant $K$.
    Then, for any $\tau\in(-\infty,1/(2\|\bfA\|K))$,
    \begin{equation*}
        \E[\exp(\tau\bZ^{\top}\bfA\bZ)]\le \exp\left(\tau\E[\bZ^{\top}\bfA\bZ]+\frac{2\|\bfA\|K\tau^{2}}{1-2\|\bfA\|K\tau}\E[\bZ^{\top}\bfA\bZ]\right).
    \end{equation*}
\end{lemma}

\begin{proof}
We divide the proof into the cases $\tau > 0$ and $\tau < 0$ (the statement with $\tau=0$ is obvious).

(Step 1: $\tau>0$)
Define $f(\bz)=\bz^{\top}\bfA\bz$ and $M_{+}(s)=\int e^{sf}\diff\nu$ for $s\ge0$ (the moment-generating function of $f(\bZ)$).
We have
\begin{equation*}
    \Ent( e^{sf})=\E[ e^{sf}sf]-\E[e^{sf}]\log\E[e^{sf}]=sM_{+}'(s)-M_{+}(s)\log M_{+}(s),
\end{equation*}
and the positive semi-definiteness of $\bfA$ yields
\begin{align*}
    \E\left[\|\nabla e^{sf/2}\|^{2}\right]&=\frac{s^{2}}{4}\int e^{sf(\bz)}\|\nabla f(\bz)\|^{2}\nu(\diff\bz)= s^{2}\int e^{sf(\bz)}\|\bfA\bz\|^{2}\nu(\diff\bz)\\
    &\le \|\bfA\|s^{2}\int e^{sf(\bz)}f(\bz)\nu(\diff\bz)= \|\bfA\|s^{2}M_{+}'(s).
\end{align*}
Since a log-Sobolev inequality holds true, we have
\begin{equation*}
    sM_{+}'(s)\le M_{+}(s)\log M_{+} (s)  +2\|\bfA\|Ks^{2}M_{+}'(s)
\end{equation*}
and thus 
\begin{equation*}
    s(1-2\|\bfA\|Ks)M_{+}'(s)\le M_{+}(s)\log M_{+}(s).
\end{equation*}
Let us define $F_{+}(s)=(1/s)\log M_{+}(s)$ (the cumulant-generating function of $f(\bZ)$ scaled by $1/s$; let $F_{+}(0):=\E[f(\bZ)]$, and then $F_{+}(s)$ is continuous since $\lim_{s\downarrow0}F_{+}(s)=\lim_{s\downarrow0}(1/s)(\log M_{+}(s)-\log M_{+}(0))=M_{+}'(0)/M_{+}(0)$) and $b:=2\|\bfA\|K$.
Then
\begin{align*}
    F_{+}'(s)&=-\frac{\log M_{+}(s)}{s^{2}}+\frac{1}{s}\frac{M_{+}'(s)}{M_{+}(s)}\\
    &\le -\frac{\log M_{+}(s)}{s^{2}}+\frac{\log M_{+}(s)}{s^{2}(1-bs)}=\left(\frac{1}{s(1-bs)}-\frac{1}{s}\right)F_{+}(s)=\left(\frac{b}{1-bs}\right)F_{+}(s).
\end{align*}
Then Gr\"{o}nwall's inequality yields
\begin{align*}
    F_{+}(s)&\le F_{+}(0)\exp\left(\int_{0}^{s}\frac{b}{1-bu}\diff u\right)=\E[f(\bZ)]\exp\left(\int_{0}^{bs}\frac{\diff v}{1-v}\right)=\E[f(\bZ)]\exp\left(\int_{1-bs}^{1}\frac{\diff v'}{v'}\right)\\
    &=\frac{\E[f(\bZ)]}{1-bs}.
\end{align*}
Therefore, we have
\begin{equation*}
    (1/s)\log M_{+}(s)\le \E[f(\bZ)]\frac{1}{1-bs},
\end{equation*}
and scaling with $s$ yields
\begin{equation*}
    \log M_{+}(s)\le s\E[f(\bZ)]\frac{1}{1-bs},
\end{equation*}
and the monotonicity of $\exp(\cdot)$ along with $1/(1-bs)=1+bs/(1-bs)$ leads to
\begin{equation*}
    M_{+}(s)\le \exp\left(s\E[f(\bZ)]\frac{1}{1-bs}\right)=\exp\left(s\E[f(\bZ)]-\E[f(\bZ)]\frac{bs^{2}}{1-bs}\right).
\end{equation*}
Hence, we have
\begin{equation*}
    \E[\exp(sf(\bZ))]\le \exp\left(s\E[f(\bZ)]+\frac{bs^{2}}{1-bs}\E[f(\bZ)]\right).
\end{equation*}

(Step 2: $\tau<0$) Let us define $g(\bz)=-\bz^{\top}\bfA\bz$ and $M_{-}(s)=\int e^{sg}\diff\nu(=\int e^{-sf}\diff\nu)$ for $s\ge0$.
Then, the positive semi-definiteness of $\bfA$ yields
\begin{align*}
    \E\left[\|\nabla e^{sg/2}\|^{2}\right]&=\frac{s^{2}}{4}\int e^{sg(\bz)}\|\nabla g(\bz)\|^{2}\nu(\diff\bz)= s^{2}\int e^{sg(\bz)}\|\bfA\bz\|^{2}\nu(\diff\bz)\\
    &\le \|\bfA\|s^{2}\int e^{sg(\bz)}|g(\bz)|\nu(\diff\bz)=\|\bfA\|s^{2}\int e^{sg(\bz)}(-g(\bz))\nu(\diff\bz)\\
    &=-\|\bfA\|s^{2}\int e^{sg(\bz)}g(\bz)\nu(\diff\bz)= -\|\bfA\|s^{2}M_{-}'(s).
\end{align*}
Therefore, a log-Sobolev inequality with constant $K$ leads to
\begin{equation*}
    sM_{-}'(s)\le M_{-}(s)\log M_{-}(s) -2\|\bfA\|Ks^{2}M_{-}'(s).
\end{equation*}
Define $F_{-}(s)=(1/s)\log M_{-}(s)$ (with $F_{-}(0)=\E[g(\bZ)]$) and $b:=2\|\bfA\|K$.
Then,
\begin{align*}
    F_{-}'(s)&=-\frac{1}{s^{2}}\log M_{-}(s)+\frac{M_{-}'(s)}{sM_{-}(s)}\\
    &\le-\frac{1}{s^{2}}\log M_{-}(s)+\frac{1}{s^{2}(1+bs)}\log M_{-}(s)=F_{-}(s)\left(\frac{1}{s(1+bs)}-\frac{1}{s}\right)=-F_{-}(s)\left(\frac{b}{1+bs}\right).
\end{align*}
Gr\"{o}nwall's inequality yields
\begin{align*}
    F_{-}(s)&\le F_{-}(0)\exp\left(-\int_{0}^{s}\frac{b}{1+bt}\diff t\right)=F_{-}(0)\exp\left(-\int_{1}^{1+bs}\frac{1}{u}\diff u\right)=F_{-}(0)\exp(-\log(1+bs))\\
    &=F_{-}(0)\frac{1}{1+bs}=\E[g(\bZ)]\frac{1}{1+bs},
\end{align*}
where the last identity is by $F_{-}(0)=\E[g]$.
Hence, we have
\begin{equation*}
    (1/s)\log M_{-}(s)\le \E[g(\bZ)]\frac{1}{1+bs},
\end{equation*}
and multiplying both sides by $s$ yields
\begin{equation*}
    \log M_{-}(s)\le s\E[g(\bZ)]\frac{1}{1+bs}.
\end{equation*}
The monotonicity of $\exp(\cdot)$ yields
\begin{equation*}
    M_{-}(s)\le \exp\left(s\E[g(\bZ)]\frac{1}{1+bs}\right)=\exp\left(s\E[g(\bZ)]-\E[g(\bZ)]\frac{bs^{2}}{1+bs}\right).
\end{equation*}
Then, we obtain
\begin{equation*}
    \E\left[\exp\left(sg(\bZ)\right)\right]\le \exp\left(s\E[g(\bZ)]-\E[g(\bZ)]\frac{bs^{2}}{1+bs}\right).
\end{equation*}
As we set $\tau=-s$, we derive the statement.
\end{proof}

\end{document}
